$[-1]^l$). In order to construct the most general pulsation, one must: (1) combine these
individual harmonic pulsations with those for $M \neq 0$, which are obtained from these by
rotations about the center of the star; (2) transform the resultant harmonic pulsations
for each value of $l$ to some common gauge (recall that the Regge-Wheeler gauge depends
upon $l$); (3) take linear combinations of the transformed harmonic pulsations.

\subsubsection*{i) \textit{Quasi-normal Modes}}

Of considerable interest are the ``quasi-normal modes'' of pulsation which, far behind
the wave front, are constructed from one single outgoing complex normal mode
\begin{equation}
  K_n(r,t) = K_n^{(0)}(r)\, e^{i\omega_n t} + K_n^{(0)*}(r)\, e^{-i\omega_n^* t}
  = \bigl[ K_n^{(0)}(r)\, e^{i\sigma_n t} + K_n^{(0)*}(r)\, e^{-i\sigma_n t} \bigr] e^{-t/\tau_n}
\label{eq:22}
\end{equation}
and similarly for $H_0, W, V$.

In the wave zone, but still far behind the wave front, the gravitational waves from such
a quasi-normal pulsation have the form (cf.\ eq.\ [15]):
\begin{align}
  K(r,t) &= 2 |C_n^{(0)}| \exp\bigl[-(t - r - 2M\ln r)/\tau_n\bigr]
             \cos\bigl[\sigma_n(t - r - 2M\ln r) + \delta_n\bigr] \notag \\
  H(r,t) &= -2\sigma_n |C_n^{(0)}| r \exp\bigl[-(t - r - 2M\ln r)/\tau_n\bigr]
             \sin\bigl[\sigma_n(t - r - 2M\ln r) + \delta_n\bigr]
             \quad \text{for } r \gg M \text{ and } r \gg 1/\sigma_n \;.
\label{eq:23}
\end{align}
Here $\delta_n$ is the phase of $C_n^{(0)}$.

Let us examine some of the properties of the quasi-normal modes (22) and (23).

\subsubsection*{ii) \textit{Stability of the Quasi-normal Modes}}

It is clear from equations (22) and (23) that the imaginary part, $1/\tau_n$, of the complex
eigenfrequency of the outgoing complex normal mode becomes, for the physical
quasi-normal mode, the damping rate or growth rate of pulsations. If $\tau_n$ is positive, the
quasi-normal pulsations are damped by the emission of gravitational waves. But if $\tau_n$
is negative, the quasi-normal ``pulsations'' are unstable, and the perturbation grows
exponentially in time. Stability versus instability is thus determined by where in the
complex plane $\omega_n$ lies: The region below the real axis is unstable; the region above the
real axis is stable.

It is actually more useful to restate this criterion in terms of $\omega_n{}^2$---the quantity appearing in the eigenequations (14). Since we have adopted the convention that $\sigma_n \geq 0$
(eq.\ [13]), there is a one-to-one relation between $\omega_n$ and $\omega_n{}^2$, and there is no ambiguity
about stability in the complex $\omega^2$ plane: \textit{For any complex outgoing normal mode} ($C_n^{(0)} \neq 0$),
\textit{the corresponding quasi-normal pulsation is stable if $\omega_n{}^2$ lies on the positive real axis
or in the upper half of the complex plane; it is unstable if $\omega_n{}^2$ lies on the negative real axis
or in the lower half of the complex plane.}

\subsubsection*{iii) \textit{Gravitational Radiation from the Quasi-normal Modes}}

Turn now from stability to the form of the gravitational waves of the quasi-normal
pulsations. From equation (23) it is clear that the gravitational radiation emitted by
the $n$th quasi-normal mode has, as seen by a man at radius $r$, the frequency
\begin{equation}
  f(r) = \frac{\sigma_n}{2\pi\,(1 - 2M/r)^{1/2}}
         \quad \text{as} \quad r \to \infty \;.
\label{eq:24}
\end{equation}
The same gravitational redshift is present here as for $r \to \infty$. Note that the logarithmic
term in the sinusoidal part of the gravitational waves (23) is unrelated to the gravita-
